\scanpage{39}
% L1–8
\uline{Théorème} On a les implications marquées sur le diagramme des notations. De plus, si on remplace dans (i) et (i bis) la relation $\operatorname{Ext}^{n-1}(k,M)=0$ par : $M$ est Coh. Mac. et si dans (v) on suppose de plus $M$ Coh. Mac., alors \deleted{toutes} les conditions \deleted{(i)} (i bis) (v) deviennent équivalentes à : $A$ est Coh. Mac. et $M\simeq\Omega_A$ (on peut alors oublier ii iii iv, qui pour $M$ Coh. Mac sont des rédites de (i) et (i bis)

% L9–11
$M$ est équidim, donc (i) $\Leftrightarrow$ (i bis). Prouvons par récurrence que alors $A$ est Coh. Mac. et $M\simeq\Omega_A$. Pour $n=0$ c'est trivial [$H^0(M)=M$ est dualisant].

% L12–22
Supposons $n>0$, et le th. prouvé pour $n'<n$. \deleted{\unsure{les}} Comme $H^n(M)$ est fidèle, \deleted{les idéaux pr. associés à $0$ de $A$ sont tous de dim. $n$, donc \unsure{et à}} \inserted{$M$ est fidèle, et étant} Coh. Mac, il s'ensuit que les idéaux pr. associés à $M$ et à $A$ sont les $=$s, et qu'ils sont tous de $=$ \deleted{dimension} \inserted{corang} $n$. Comme $n>0$, il y a un $x$ qui n'est contenu dans aucun d'eux, donc qui n'est pas div. de $0$ dans $A$ ou dans $M$. \deleted{\unsure{Donc on|D'après}} Ecrivons alors

% L23–26
\[
 \operatorname{Ext}_{A}^{\unsure{*|\bullet}}(N,M)\Longleftarrow
 \operatorname{Ext}_{A/x}^{*}(N,\operatorname{Ext}_{A}^{*}(A/x,M))
\]
On a $\operatorname{Ext}_{A}^{0}(A/x,M)=0$, $\operatorname{Ext}_{A}^{1}(A/x,M)=M/xM$, d'où
\[
 \operatorname{Ext}_{A}^{i}(N,M)\simeq\operatorname{Ext}_{A/x}^{\unsure{i-1|i1}}(N,M/xM)
\]
pour tout $N$. Par suite, la condition \unsure{(i a)|(i'a)}

% archive sheet number: 38
